\documentclass[a4paper,11pt]{article}
\usepackage[T1]{fontenc}
\usepackage[utf8]{inputenc}
\usepackage[french]{babel}
\usepackage{geometry}
\geometry{top=1.5cm, bottom=1.8cm, left=1.8cm, right=1.8cm}
\usepackage{xcolor}
\renewcommand{\familydefault}{\sfdefault}
\definecolor{Couleur1}{HTML}{8b1f30}
\definecolor{Couleur6}{HTML}{CF6876}
\definecolor{Couleur7}{HTML}{A5C03F}
\definecolor{Couleur8}{HTML}{D36740}
\usepackage{amsmath, amssymb}
\usepackage{array, colortbl}
\usepackage{tikz}
\usepackage{fontawesome5}
\usepackage{mdframed}
\usepackage{fancyhdr}
\usepackage{enumitem}
\usepackage{multicol}
\usepackage{rotating}
\setlength{\parindent}{0pt}
\newcommand{\bl}[1][4cm]{\,\makebox[#1]{\dotfill}\,}
\newcommand{\lignepoints}{\noindent\dotfill\vspace{0.3em}\par}
\newcommand{\coeur}{\textcolor{Couleur8}{\faHeart}\ }
\newcommand{\coeurrappel}{\textcolor{Couleur6}{\faHeart}\ }
\pagestyle{fancy}
\fancyhf{}
\renewcommand{\headrulewidth}{0pt}
\fancyfoot[L]{\small Mme Bertail}
\fancyfoot[C]{\small \textcolor{Couleur6}{Interro de le\c{c}on n\degre~15 -- Ch 1 a 19}}
\fancyfoot[R]{\small \thepage}
\begin{document}
\begin{mdframed}[linecolor=Couleur1,linewidth=2pt,topline=true,bottomline=true,leftline=false,rightline=false,innertopmargin=8pt,innerbottommargin=8pt,backgroundcolor=Couleur1!5]
\begin{center}
{\Large\bfseries\textcolor{Couleur1}{Interrogation de le\c{c}on n\degre~15}}\\[4pt]
{\large\textcolor{Couleur8}{Chapitre 19 --- Probabilit\'es}}
\end{center}
\end{mdframed}
\vspace{0.4cm}
\noindent\textbf{Nom :} \dotfill\quad\textbf{Pr\'enom :} \dotfill\quad\textbf{Classe :} \dotfill\quad\textbf{Date :} \dotfill
\vspace{0.5cm}

\vspace{0.3cm}
\noindent{\textcolor{Couleur8}{\rule{\textwidth}{0.5pt}}}\\
{\small\textit{\textcolor{Couleur8}{Nouveau --- Chapitre 19}}}
\vspace{0.1cm}

\noindent \coeur \textbf{Compl\'eter les d\'efinitions.}
\begin{mdframed}[linecolor=Couleur6,linewidth=4pt,topline=false,bottomline=false,rightline=false,backgroundcolor=Couleur6!10,innerleftmargin=8pt,innertopmargin=6pt,innerbottommargin=6pt]
\begin{itemize}[label=\textcolor{Couleur6}{$\bullet$},nosep,itemsep=4pt]
\item Une exp\'erience est \textbf{al\'eatoire} lorsque \bl[7cm].
\item Une \textbf{issue} est \bl[7cm].
\item Un \textbf{\'ev\'enement} d'une exp\'erience al\'eatoire est \bl[6cm].
\end{itemize}
\end{mdframed}

\noindent \coeur \textbf{Compl\'eter la d\'efinition.}
\begin{mdframed}[linecolor=Couleur6,linewidth=4pt,topline=false,bottomline=false,rightline=false,backgroundcolor=Couleur6!10,innerleftmargin=8pt,innertopmargin=6pt,innerbottommargin=6pt]
La \textbf{probabilit\'e} d'un \'ev\'enement donne une mesure de la chance que cet \'ev\'enement \bl[4cm].\\[4pt]
La probabilit\'e d'un \'ev\'enement est un nombre compris entre \bl[2cm] et \bl[2cm].
\end{mdframed}

\vspace{0.3cm}
\noindent{\textcolor{Couleur8}{\rule{\textwidth}{0.5pt}}}\\
{\small\textit{\textcolor{Couleur8}{Rappels --- Chapitres pr\'ec\'edents}}}
\vspace{0.1cm}

\noindent \coeurrappel Le \textbf{p\'erim\`etre} d'un polygone est \dotfill\\[4pt]
$P_{\text{Rectangle}} = $ \bl[4cm] \quad ; \quad $P_{\text{Carr\'e}} = $ \bl[3cm] \quad ; \quad $P_{\text{Disque}} = $ \bl[4cm]

\noindent \coeurrappel L'\textbf{aire} d'une figure est \dotfill\\[4pt]
$A_{\text{Rectangle}} = $ \bl[4cm] \quad ; \quad $A_{\text{Carr\'e}} = $ \bl[4cm]
\end{document}